\documentclass{fiche}
\metadonnees{17}{Une vérité universelle}{Bloc 4 — Hypothèse, parole rapportée, traduction}

\begin{document}
\entete

\objectifs{%
\textbf{Langue} : discours indirect ; concordance des temps ; verbes introducteurs.\par
\textbf{Texte} : Jane Austen, \emph{Pride and Prejudice} (1813), chapitre \textsc{i}.\par
\textbf{Durée} : 1\,h\,30 — exercices 35 min, texte et questions 30 min, version 25 min.}

%% ===================================================================
\exercice[14 min]{Passer au discours indirect}

\rappel{\textbf{Rappel.} Quand le verbe introducteur est au passé, tous les temps reculent
d'un cran, et les repères de temps et de lieu changent avec le point de vue.}

\consigne{Rapportez ces paroles.}

\begin{anglais}
\begin{enumerate}[label=\arabic*.,itemsep=2.5mm,leftmargin=*]
  \item ``Netherfield Park is let at last.'' \emph{(She told him\ldots)}
    \lignes{1}
    \corr{She told him that Netherfield Park was let at last.}
  \item ``Mrs Long has just been here.'' \emph{(She said\ldots)}
    \lignes{1}
    \corr{She said that Mrs Long had just been there.}
  \item ``I have no objection to hearing it.'' \emph{(He replied\ldots)}
    \lignes{1}
    \corr{He replied that he had no objection to hearing it.}
  \item ``What is his name?'' \emph{(He asked\ldots)}
    \lignes{1}
    \corr{He asked what his name was.}
  \item ``Is he married or single?'' \emph{(He asked\ldots)}
    \lignes{1}
    \corr{He asked whether he was married or single.}
  \item ``You must visit him as soon as he comes.'' \emph{(She insisted that\ldots)}
    \lignes{1}
    \corr{She insisted that he must (\emph{ou} had to) visit him as soon as he came.}
  \item ``I will send a few lines by you.'' \emph{(He promised\ldots)}
    \lignes{1}
    \corr{He promised that he would send a few lines by her.}
  \item ``Do not talk of your nerves again.'' \emph{(He begged her\ldots)}
    \lignes{1}
    \corr{He begged her not to talk of her nerves again.}
  \item ``I saw him yesterday.'' \emph{(She claimed\ldots)}
    \lignes{1}
    \corr{She claimed that she had seen him the day before.}
  \item ``We shall meet him at the assemblies tomorrow.'' \emph{(Elizabeth reminded her
    mother\ldots)}
    \lignes{1}
    \corr{Elizabeth reminded her mother that they would meet him at the assemblies the next
    day.}
\end{enumerate}
\end{anglais}

\note[Le tableau des reculs]{présent $\rightarrow$ prétérit ; present perfect et prétérit
$\rightarrow$ pluperfect ; \emph{will} $\rightarrow$ \emph{would} ; \emph{can}
$\rightarrow$ \emph{could} ; \emph{may} $\rightarrow$ \emph{might}. \emph{Must},
\emph{should}, \emph{ought to} et \emph{would} ne reculent pas : ils n'ont pas de forme
antérieure.}

\note[Les repères aussi]{\emph{here} $\rightarrow$ \emph{there}, \emph{this} $\rightarrow$
\emph{that}, \emph{now} $\rightarrow$ \emph{then}, \emph{today} $\rightarrow$ \emph{that
day}, \emph{yesterday} $\rightarrow$ \emph{the day before}, \emph{tomorrow} $\rightarrow$
\emph{the next day}, \emph{ago} $\rightarrow$ \emph{before} (fiche 5).}

%% ===================================================================
\exercice[10 min]{Le verbe introducteur}

\rappel{\textbf{Rappel.} \emph{Say} n'a pas de complément de personne, \emph{tell} en exige
un. Les autres verbes introducteurs disent en un mot ce que fait celui qui parle.}

\consigne{Complétez avec le verbe qui convient : \emph{said}, \emph{told}, \emph{asked},
\emph{replied}, \emph{admitted}, \emph{suggested}, \emph{complained}, \emph{warned},
\emph{promised}, \emph{denied}.}

\begin{anglais}
\begin{enumerate}[label=\arabic*.,itemsep=2.2mm,leftmargin=*]
  \item She \trou{told} him that Netherfield was let.
  \item She \trou{said} that Mrs Long had told her everything.
  \item He \trou{asked} her what the young man's name was.
  \item He \trou{replied} that he had no objection to hearing it.
  \item She \trou{complained} that he had no compassion on her nerves.
  \item He \trou{suggested} sending the girls by themselves.
  \item She \trou{denied} having said any such thing.
  \item He \trou{promised} to send a few lines by her.
  \item She \trou{admitted} that she had had her share of beauty.
  \item He \trou{warned} her that twenty such men would be of no use.
\end{enumerate}
\end{anglais}

\note[Trois constructions]{\emph{Suggest} et \emph{deny} sont suivis d'un gérondif, non d'un
infinitif : \emph{he suggested sending}, \emph{she denied having said} --- c'est le sujet de
la fiche 18. \emph{Promise}, \emph{offer}, \emph{refuse} prennent l'infinitif.
\emph{Warn}, \emph{remind}, \emph{tell} exigent un complément de personne.}

%% ===================================================================
\exercice[11 min]{Le mariage et la fortune}
\consigne{a) Associez chaque mot à sa traduction.}

\begin{multicols}{2}
\begin{enumerate}[label=\arabic*.,itemsep=1.2mm,leftmargin=*]
  \item to be let \hfill \trou{f}
  \item an estate \hfill \trou{c}
  \item an establishment \hfill \trou{k}
  \item a fortune \hfill \trou{a}
  \item to settle \hfill \trou{h}
  \item a design \hfill \trou{b}
  \item to flatter \hfill \trou{j}
  \item tiresome \hfill \trou{d}
  \item to vex \hfill \trou{i}
  \item mean \hfill \trou{e}
  \item a solace \hfill \trou{g}
  \item to give over \hfill \trou{l}
\end{enumerate}
\columnbreak
\begin{enumerate}[label=\alph*.,itemsep=1.2mm,leftmargin=*]
  \item une fortune
  \item un dessein, une intention
  \item un domaine \textit{(fiche 8)}
  \item assommant, agaçant
  \item médiocre, borné
  \item être loué (maison)
  \item une consolation
  \item s'installer, se fixer
  \item contrarier, agacer
  \item flatter
  \item un établissement, une situation
  \item cesser de
\end{enumerate}
\end{multicols}

\note[Deux faux amis d'époque]{\emph{An establishment} désigne ici la situation qu'un mariage
procure à une jeune fille, non un établissement scolaire ou commercial. \emph{Mean} appliqué à
l'esprit signifie \og borné \fg, comme dans \emph{a woman of mean understanding} : c'est le
même mot que \emph{the meanest reptile} de la fiche 7.}

\consigne{b) Complétez avec un mot de la liste.}
\begin{anglais}
\begin{enumerate}[label=\arabic*.,itemsep=2mm,leftmargin=*]
  \item Netherfield Park \trou{is let} at last, to a young man of large \trou{fortune}.
  \item ``Is that his \trou{design} in \trou{settling} here?'' asked Mr Bennet.
  \item ``You take delight in \trou{vexing} me,'' she answered.
  \item When a woman has five grown-up daughters, she ought to \trou{give over} thinking of
    her own beauty.
\end{enumerate}
\end{anglais}

%% ===================================================================
\newpage
\section{Texte}

\noindent\small\textit{Les Bennet ont cinq filles et peu de fortune. Une grande propriété du
voisinage vient d'être louée.}\normalsize

\begin{texte}
``My dear Mr. Bennet,'' said his lady to him one day, ``have you heard that Netherfield Park
is \gl[to let]{let}{loué} at last?''

Mr. Bennet replied that he had not.

``But it is,'' returned she; ``for Mrs. Long has just been here, and she told me all about
it.''

Mr. Bennet made no answer.

``Do not you want to know who has taken it?'' cried his wife, impatiently.

``\emph{You} want to tell me, and I have no objection to hearing it.''

This was invitation enough.

``Why, my dear, you must know, Mrs. Long says that Netherfield is taken by a young man of
large fortune from the north of England; that he came down on Monday in a
\gl[a chaise and four]{chaise and four}{une voiture à quatre chevaux} to see the place, and
was so much delighted with it that he agreed with Mr. Morris immediately; that he is to take
possession before \gl[Michaelmas]{Michaelmas}{la Saint-Michel, le 29 septembre}, and some of
his servants are to be in the house by the end of next week.''

``What is his name?''

``Bingley.''

``Is he married or single?''

``Oh, single, my dear, to be sure! A single man of large fortune; four or five thousand a
year. What a fine thing for our girls!''

``How so? how can it affect them?''

``My dear Mr. Bennet,'' replied his wife, ``how can you be so
\gl[tiresome]{tiresome}{assommant}? You must know that I am thinking of his marrying one of
them.''

``Is that his \gl[a design]{design}{un dessein, une intention} in settling here?''

``Design? Nonsense, how can you talk so! But it is very likely that he \emph{may} fall in love
with one of them, and therefore you must visit him as soon as he comes.''

``I see no occasion for that. You and the girls may go --- or you may send them by themselves,
which perhaps will be still better; for as you are as handsome as any of them, Mr. Bingley
might like you the best of the party.''

``My dear, you flatter me. I certainly \emph{have} had my share of beauty, but I do not
pretend to be anything extraordinary now. When a woman has five grown-up daughters, she ought
to \gl[to give over]{give over}{cesser de} thinking of her own beauty.''

``In such cases, a woman has not often much beauty to think of.''

``But, my dear, you must indeed go and see Mr. Bingley when he comes into the neighbourhood.''

``It is more than I \gl[to engage for]{engage for}{s'engager à} , I assure you.''

``But consider your daughters. Only think what an
\gl[an establishment]{establishment}{une situation, un établissement} it would be for one of
them. Sir William and Lady Lucas are determined to go, merely on that account; for in general,
you know, they visit no new comers. Indeed you must go, for it will be impossible for
\emph{us} to visit him, if you do not.''

``You are over \gl[scrupulous]{scrupulous}{scrupuleux}, surely. I dare say Mr. Bingley will be
very glad to see you; and I will send a few lines by you to assure him of my hearty consent to
his marrying whichever he chooses of the girls --- though I must throw in a good word for my
little Lizzy.''

``I desire you will do no such thing. Lizzy is not a bit better than the others: and I am sure
she is not half so handsome as Jane, nor half so good-humoured as Lydia. But you are always
giving \emph{her} the preference.''

``They have none of them much to recommend them,'' replied he: ``they are all silly and
ignorant like other girls; but Lizzy has something more of quickness than her sisters.''

``Mr. Bennet, how can you abuse your own children in such a way? You take delight in
\gl[to vex]{vexing}{contrarier, agacer} me. You have no compassion on my poor nerves.''

``You mistake me, my dear. I have a high respect for your nerves. They are my old friends. I
have heard you mention them with consideration these twenty years at least.''
\end{texte}
\source{Jane Austen, \emph{Pride and Prejudice}, Londres, Egerton, 1813, ch.~\textsc{i}.}

\partiel[15 min]{Questions}
\consigne{\en{Answer in English, in full sentences.}}

\begin{enumerate}[label=\arabic*.,itemsep=2mm,leftmargin=*]
  \item \en{What news does Mrs Bennet bring, and from whom does she have it?}
    \lignes{2}
    \corr{That Netherfield Park has been let to a young man of large fortune from the north
    of England. She has it from Mrs Long, who has just called on her.}
  \item \en{``You want to tell me, and I have no objection to hearing it.'' What does this
    answer show about Mr Bennet?}
    \lignes{3}
    \corr{That he has seen through her: she is not asking a question, she is waiting to be
    asked. He refuses to play the part she wants, and gives her the information she wanted
    anyway --- which is his whole manner in the scene.}
  \item \en{Why does Mrs Bennet want her husband to call on Mr Bingley?}
    \lignes{3}
    \corr{Because in that society a family cannot visit a newcomer until the father has
    called first. Without his visit, her daughters cannot meet him, and Lady Lucas's
    daughters will get there first.}
  \item \en{How does Mr Bennet answer when she says he flatters her?}
    \lignes{2}
    \corr{\emph{In such cases, a woman has not often much beauty to think of} --- which
    withdraws the compliment he had just paid her, without her noticing.}
  \item \en{What does each parent think of Lizzy?}
    \lignes{3}
    \corr{Her father says all five girls are silly and ignorant but that Lizzy has something
    more of quickness, and he throws in a good word for her. Her mother says she is not a bit
    better than the others, less handsome than Jane and less good-humoured than Lydia.}
  \item \en{``I have a high respect for your nerves. They are my old friends.'' Explain the
    joke.}
    \lignes{3}
    \corr{Her nerves are not an illness but a habit: he has been hearing about them for
    twenty years, so they are as familiar to him as old acquaintances. The politeness of the
    sentence is what makes it cruel.}
\end{enumerate}

\note[Discours direct et indirect dans le texte]{\emph{Mr. Bennet replied that he had not} :
le seul discours indirect du chapitre, glissé au milieu des répliques --- il résume ce qui
n'a aucun intérêt. Et dans la longue tirade de Mrs Bennet, \emph{Mrs Long says that\ldots{}
that\ldots{} that\ldots} : trois subordonnées enchaînées, où l'on entend que tout est de
seconde main.}

%% ===================================================================
\newpage
\section{Version}
\consigne{Traduisez le passage en français. Il s'agit des deux premiers paragraphes du roman,
du dernier paragraphe du chapitre, et de l'ouverture du chapitre suivant.}

\begin{version}
It is a truth universally acknowledged, that a single man in possession of a good fortune must
be in want of a wife.

However little known the feelings or views of such a man may be on his first entering a
neighbourhood, this truth is so well fixed in the minds of the surrounding families, that he
is considered as the rightful property of some one or other of their daughters.

[\ldots]

Mr. Bennet was so odd a mixture of quick \gl[parts]{parts}{les dons, l'esprit}, sarcastic
humour, reserve, and \gl[caprice]{caprice}{la fantaisie}, that the experience of
three-and-twenty years had been insufficient to make his wife understand his character.
\emph{Her} mind was less difficult to \gl[to develope]{develope}{sonder, développer ---
orthographe ancienne de} develop. She was a woman of \gl[mean]{mean}{borné, médiocre}
understanding, little information, and uncertain temper. When she was discontented, she
fancied herself nervous. The business of her life was to get her daughters married: its
\gl[a solace]{solace}{une consolation} was visiting and news.

[\ldots]

Mr. Bennet was among the earliest of those who \gl[to wait on]{waited on}{rendre visite à,
faire sa cour à} Mr. Bingley. He had always intended to visit him, though to the last always
assuring his wife that he should not go; and till the evening after the visit was paid she had
no knowledge of it.
\end{version}
\source{\emph{Ibid.}, ch.~\textsc{i} et \textsc{ii}.}

\lignes{20}

\corr{\textbf{Proposition de traduction.} C'est une vérité universellement reconnue qu'un
célibataire pourvu d'une belle fortune doit avoir besoin d'une épouse.

Si peu que l'on connaisse les sentiments ou les vues d'un tel homme lorsqu'il arrive dans une
région, cette vérité est si bien ancrée dans l'esprit des familles d'alentour qu'on le
considère comme la propriété légitime de l'une ou l'autre de leurs filles.

[\ldots]

M. Bennet était un si curieux mélange de vivacité d'esprit, d'humour sarcastique, de réserve
et de fantaisie, que l'expérience de vingt-trois ans n'avait pas suffi à sa femme pour
comprendre son caractère. L'esprit de celle-ci était moins difficile à sonder. C'était une
femme d'intelligence bornée, de peu d'instruction et d'humeur inégale. Quand elle était
mécontente, elle se croyait nerveuse. L'affaire de sa vie était de marier ses filles ; sa
consolation, les visites et les nouvelles.

[\ldots]

M. Bennet fut des tout premiers à rendre visite à M. Bingley. Il avait toujours eu
l'intention d'y aller, tout en assurant sa femme jusqu'au bout qu'il n'irait pas ; et ce ne
fut que le soir du lendemain de sa visite qu'elle l'apprit.}

\note[Choix de traduction 1]{\emph{must be in want of a wife} : \emph{must} de déduction
(fiche 8), non d'obligation, et \emph{to be in want of} = \og avoir besoin de \fg. La phrase
la plus célèbre du roman se joue sur ce \emph{must} : ce n'est pas une nécessité, c'est ce
que tout le voisinage décide de croire.}

\note[Choix de traduction 2]{\emph{However little known\ldots may be} : concession avec
\emph{however} + adjectif en tête. Le français dit \og si peu que \fg{} + subjonctif, ou
\og quelque peu connus que soient \fg. \emph{May} y marque la concession, non la
possibilité.}

\note[Choix de traduction 3]{\emph{so odd a mixture of\ldots that} : l'ordre \emph{so +
adjectif + a + nom} est obligatoire après \emph{so} ; le français rétablit \og un si curieux
mélange \fg.}

\note[Choix de traduction 4]{\emph{always assuring his wife that he should not go} : \emph{he
should not go} est un discours indirect --- au style direct, \emph{I shall not go}.
\emph{Shall} devient \emph{should} à la première personne comme \emph{will} devient
\emph{would} ailleurs. Usage d'époque, à reconnaître.}

\note[Choix de traduction 5]{\emph{till the evening after the visit was paid she had no
knowledge of it} : le français préfère retourner la négation en \og ce ne fut que\ldots
qu'elle l'apprit \fg. Traduire \emph{not\ldots till} par \og pas avant \fg{} est correct mais
plus plat.}

%% ===================================================================
\partiel{Lexique à mémoriser}
\begin{itemize}[label=--,itemsep=0.8mm,leftmargin=*]\small
  \item \textbf{to let} — louer (donner en location) ; to rent \textit{« louer, prendre en location »}
  \item \textbf{a fortune} — une fortune ; a man of large fortune
  \item \textbf{to settle} — s'installer ; \textit{noms} a settlement, a settler
  \item \textbf{a design} — un dessein ; \textit{et} un dessin, un modèle
  \item \textbf{an establishment} — une situation (par mariage)
  \item \textbf{a neighbourhood} — le voisinage ; neighbour + hood\textit{, suffixe collectif comme dans} childhood
  \item \textbf{acknowledged} — reconnu ; \textit{verbe} to acknowledge
  \item \textbf{rightful} — légitime ; right + -ful
  \item \textbf{tiresome} — assommant ; tire + -some\textit{, comme} troublesome, quarrelsome
  \item \textbf{to vex} — contrarier
  \item \textbf{to abuse} — dire du mal de ; \textit{et} maltraiter --- \textit{faux ami :} to deceive \textit{pour « abuser »}
  \item \textbf{mean} — borné, mesquin \textit{(fiche 7)}
  \item \textbf{quickness} — la vivacité d'esprit ; quick + -ness
  \item \textbf{a solace} — une consolation
  \item \textbf{to fancy} — s'imaginer ; \textit{nom} a fancy \textit{« une lubie »}
  \item \textbf{handsome} — beau ; \textit{s'appliquait alors aux deux sexes}
  \item \textbf{good-humoured} — d'humeur agréable
  \item \textbf{to dare say} — supposer ; I dare say \textit{« j'imagine »}
  \item \textbf{to take delight in} — prendre plaisir à
  \item \textbf{hearty} — chaleureux, cordial ; heart + -y
\end{itemize}

\end{document}
